\documentclass[11pt]{article}
\usepackage[UTF8]{ctex}
\usepackage{amsmath,amssymb,amsthm}
\usepackage{etoolbox}

% 1. 自定义定理样式实现标题后强制换行
\newtheoremstyle{break}% 样式名称
  {\topsep}{\topsep}% 环境上下间距
  {\normalfont}% 正文字体
  {}% 缩进量
  {\bfseries}{.}% 头部字体和标点
  {\newline}% 头部后的间距：设置为 \newline 即可实现完美的换行效果
  {}% 头部时间/附加参数规范

% 应用新样式并声明常用的定理类环境
\theoremstyle{break}
\newtheorem{theorem}{定理}[section]
\newtheorem{lemma}[theorem]{引理}
\newtheorem{corollary}[theorem]{推论}
\newtheorem{definition}[theorem]{定义}

% 2. 利用 etoolbox 宏包动态修补 proof 环境，使其标题后自动换行
\patchcmd{\proof}{\item[\hskip\labelsep\itshape#1\@addpunct{.}]}{\item[\hskip\labelsep\itshape#1\@addpunct{.}]\hfill\par}{}{}

\begin{document}
\begin{center}
\textbf{固定重数和集大小的多项式压缩} \\
\small 草稿 \quad 2026年4月29日
\end{center}
\noindent\textbf{摘要} \\
对于固定的 $h$，令 $N(h,k)$ 为满足如下条件的最小正整数 $N$：任意 $k$ 元整数集所能取到的 $h$ 重和集基数 $|hA|$，均可由 $\{1,\dots,N\}$ 的某个 $k$ 元子集取到。利用 Rajagopal 关于 $|hA|$ 可能取值范围的固定 $h$ 定理，我们证明了对任意固定的 $h$，均有 $N(h,k) \le k^{C_h}$。本文的核心创新在于，以多项式直径的构造替换了 Rajagopal 大值构造中的稀疏几何块。该替换操作在低次 Hilbert 函数的层面上完成，因为当 $h$ 固定时，Hilbert 函数是刻画 $h$ 重和集的唯一相关数据。
\section{引言}
设 $A \subset \mathbb{Z}$ 为有限集，$h$ 为正整数，记
\[
hA = \{a_1 + \cdots + a_h : a_i \in A\},
\]
其中允许元素重复。定义
\[
R(h,k) = \{|hA| : A \subset \mathbb{Z}, |A| = k\}.
\]
令 $N(h,k)$ 为满足
\[
R(h,k) = \{|hA| : A \subseteq \{1,\dots,N\}, |A| = k\}
\]
的最小正整数 $N$。等价地，我们可以在区间 $[0, N-1]$ 上进行讨论，并在最后进行平移。
本文旨在证明如下多项式压缩界。
\textbf{定理 1.1.}
对任意固定的 $h \ge 1$，存在常数 $C$，使得对所有 $k \ge 2$，均有
\[ N(h,k) \le k^{Ch}. \]
证明利用了 Rajagopal 关于可能的固定重数求和集大小的定理。令
\[ M = \binom{h+k-1}{h}, \quad K = \binom{h+k-1}{k}, \]
并定义
\[ \Delta = \bigcup_{\ell=0}^{\min(h,k)-3} \left[ M + \ell h + 1, \, M + \ell h + (h-2-\ell) \right]. \]
此处及下文中，$[u,v]$ 表示满足 $u \le n \le v$ 的整数 $n$ 构成的集合；若 $u > v$，则该区间为空集。
\end{document}